\documentclass[11pt]{article}

\usepackage[margin=1in]{geometry}
\usepackage{amsmath,amssymb,mathtools}
\usepackage{microtype}
\usepackage{newtxtext,newtxmath}
\usepackage{enumitem}

\setlist{nosep}
\setlength{\parindent}{1.2em}
\setlength{\parskip}{0.35em}

\newcommand{\Aspace}{\mathcal{A}}
\DeclareMathOperator{\sech}{sech}
\DeclareMathOperator{\csch}{csch}
\DeclareMathOperator{\gd}{gd}
\DeclareMathOperator{\artanh}{artanh}

\title{\texttt{sin()} Stays \texttt{sin()}; Angles Become Plural\
\large A Breathing Pass on Circular, Hyperbolic, and Reciprocal Vantage Coordinates}
\author{Wilder Blair Munro $\times$ Aurora}
\date{Version 2.5 --- expanded working draft}

\begin{document}
\maketitle

\begin{abstract}
For a one-dimensional relativistic speed parameter $\beta=v/c$ with
$0<\beta<1$, one may write

$$
\beta=\sin\theta=\tanh\chi,
\qquad
\gamma=\sec\theta=\cosh\chi,
$$

where $\chi=\artanh\beta$ is rapidity and $\theta=\gd\chi$ is its
Gudermannian circular image. If one also introduces reciprocal coordinates

$$
\alpha=\beta^{-1},
\qquad
\delta=(1-\beta^2)^{-1/2},
$$

then

$$
\alpha=\csc\theta=\coth\chi,
\qquad
\delta=\sec\theta=\cosh\chi.
$$

These identities are standard consequences of elementary trigonometry,
hyperbolic parametrization, and the Gudermannian. The proposal examined here
is not a new trigonometric function and not a modification of special
relativity. It is a change in what is promoted to the status of an angular
object: circular angle, rapidity, and reciprocal aperture are treated as
different coordinate bodies of one underlying one-parameter vantage relation.

This expanded draft separates the established identities from that
interpretive promotion, makes the reciprocal-triangle construction explicit,
explains why the notation $\sin\alpha=1/\alpha$ is ordinarily false but
structurally suggestive, and states precisely what a proposed
lift--rotate--squeeze realization would still have to prove. Particular care
is given to the distinction between a global relativistic complement and the
componentwise quantities $c_i$ suggested by the earlier General Peace
Dynamics construction.
\end{abstract}

\section{Scope: what is mathematics, and what is the proposal?}

The shortest version of this paper is dangerously easy to misread. Three
logically distinct layers must therefore be separated before any compression
is attempted.

The first layer is ordinary special-relativistic parametrization. For a
collinear boost, rapidity $\chi$ satisfies
\begin{equation}
\beta=\tanh\chi,
\qquad
\gamma=\cosh\chi,
\qquad
\gamma\beta=\sinh\chi.
\label{eq:rapidity}
\end{equation}
Nothing proposed here alters these relations.

The second layer is ordinary mathematics connecting circular and hyperbolic
functions. The Gudermannian~\cite{DLMF} supplies a real-valued map between a
hyperbolic parameter $\chi\in\mathbb{R}$ and a circular angle
$\theta\in(-\pi/2,\pi/2)$ without introducing an imaginary argument. In
particular,
\begin{equation}
\theta=\gd\chi
\label{eq:gddef}
\end{equation}
implies
\begin{equation}
\sin\theta=\tanh\chi,
\qquad
\cos\theta=\sech\chi,
\qquad
\tan\theta=\sinh\chi.
\label{eq:gdidentities}
\end{equation}
The reciprocal identities
\begin{equation}
\csc\theta=\coth\chi,
\qquad
\sec\theta=\cosh\chi,
\qquad
\cot\theta=\csch\chi
\label{eq:gdreciprocal}
\end{equation}
follow immediately.

The third layer is the speculative step. We ask whether the several
parameters that encode the same one-dimensional relativistic state should be
regarded merely as convenient substitutions, or instead as coordinate
representations of a more primitive geometric relation. The working name for
that relation is a \emph{vantage angle}. In this language, the proposal is
not that $\alpha$, $\theta$, and $\chi$ are numerically equal. They are not.
The proposal is that they may be coordinates of the same underlying state in
different charts or representations.

This distinction is the entire paper. Every identity may remain ordinary
while the ontology of ``angle'' changes.

\section{Minimal GPD vocabulary, without requiring GPD}

The terminology motivating this paper comes from earlier General Peace
Dynamics (GPD) work~\cite{MunroGPD}, in which a \emph{vantage} is treated
broadly as a reference-frame-like description and transformations between
vantages are promoted as primary objects. Earlier GPD notes use this strategy
far beyond special relativity, including observer--individual dual
descriptions and a real seven-dimensional construction in which spatial
coordinates are paired with corresponding time-like or phase-like
coordinates. Those broader claims are not required here.

Only four pieces of vocabulary are needed for the present argument.

\begin{enumerate}[label=(\roman*)]
\item A \emph{vantage} is a choice of representation of a physical relation.
\item A \emph{lift} embeds a description into a larger state space in which a
desired relation may have a simpler geometric representation.
\item A \emph{rotation} changes orientation within that lifted space.
\item A \emph{squeeze} or projection returns the lifted description to the
lower-dimensional variables in which measurements are reported.
\end{enumerate}

The phrase \emph{lift--rotate--squeeze} therefore denotes a proposed
factorization of a change of physical description. It should not be read as a
derivation merely because it reproduces a familiar formula. A legitimate
derivation must independently motivate the lift, specify the metric or
geometric structure of the lifted space, specify the rotation, specify the
squeeze, and demonstrate that their composition has the required Lorentzian
properties.

The present paper addresses only one missing piece of that motivation: why a
circular rotation might appear at all in a theory whose conventional boost
parameter is hyperbolic.

\section{The one-axis relation}

Consider motion along one spatial axis. Restrict initially to the positive
branch
\begin{equation}
0<v<c,
\qquad
\beta:=\frac{v}{c}\in(0,1).
\label{eq:beta}
\end{equation}
Define the complementary dimensionless quantity
\begin{equation}
\beta_c:=\sqrt{1-\beta^2}.
\label{eq:betac}
\end{equation}
Then
\begin{equation}
\beta^2+\beta_c^2=1.
\label{eq:unitcircle}
\end{equation}
Equation \eqref{eq:unitcircle} is simply the unit circle. Consequently there
exists a unique
\begin{equation}
\theta\in(0,\pi/2)
\end{equation}
such that
\begin{equation}
\beta=\sin\theta,
\qquad
\beta_c=\cos\theta.
\label{eq:circularparam}
\end{equation}

At the same time, define the ordinary Lorentz factor
\begin{equation}
\gamma:=\frac{1}{\sqrt{1-\beta^2}}.
\label{eq:gamma}
\end{equation}
Using \eqref{eq:betac},
\begin{equation}
\gamma=\frac{1}{\beta_c}.
\label{eq:gammabetac}
\end{equation}
Thus the same one-parameter state can already be written in two elementary
forms:
\begin{equation}
(\beta,\gamma)
==============

(\sin\theta,\sec\theta).
\label{eq:circpair}
\end{equation}

Now introduce rapidity
\begin{equation}
\chi:=\artanh\beta.
\label{eq:chi}
\end{equation}
Then
\begin{equation}
\beta=\tanh\chi,
\qquad
\beta_c=\sech\chi,
\qquad
\gamma=\cosh\chi.
\label{eq:hyperpair}
\end{equation}
Comparing \eqref{eq:circularparam} and \eqref{eq:hyperpair} gives
\begin{equation}
\sin\theta=\tanh\chi,
\qquad
\cos\theta=\sech\chi,
\label{eq:bridge}
\end{equation}
which is exactly the Gudermannian relation \eqref{eq:gddef}.

There is therefore no need to invent a circular--hyperbolic bridge. It
already exists.

\section{The reciprocal representation}

The earlier GPD derivation introduced quantities of the form
\begin{equation}
\alpha:=\frac{1}{\beta},
\qquad
\delta:=\frac{1}{\beta_c}.
\label{eq:alphadelta}
\end{equation}
For the positive one-axis branch considered here,
\begin{equation}
\alpha>1,
\qquad
\delta>1.
\end{equation}
Using \eqref{eq:circularparam},
\begin{equation}
\alpha=\csc\theta,
\qquad
\delta=\sec\theta.
\label{eq:alphadeltacircular}
\end{equation}
Using the Gudermannian identities,
\begin{equation}
\alpha=\coth\chi,
\qquad
\delta=\cosh\chi.
\label{eq:alphadeltahyper}
\end{equation}
Thus the reciprocal quantities are not arbitrary decorations placed on the
circular triangle. They are standard reciprocal trigonometric and hyperbolic
functions evaluated on the same underlying one-parameter state.

The constraint \eqref{eq:unitcircle} becomes
\begin{equation}
\frac{1}{\alpha^2}+\frac{1}{\delta^2}=1.
\label{eq:reciprocalconstraint}
\end{equation}
Multiplying by $\alpha^2\delta^2$ gives
\begin{equation}
\alpha^2+\delta^2=\alpha^2\delta^2.
\label{eq:reciprocalpythagoras}
\end{equation}
Since $\alpha$ and $\delta$ are positive,
\begin{equation}
\sqrt{\alpha^2+\delta^2}=\alpha\delta.
\label{eq:reciprocalhypotenuse}
\end{equation}
Equation \eqref{eq:reciprocalhypotenuse} is the exact algebraic basis of the
reciprocal triangle that appeared in the earlier sketch work.

\section{The reciprocal triangle is the same triangle at another scale}

The normalized circular triangle has horizontal leg $\beta_c$, vertical leg
$\beta$, and hypotenuse $1$. Write it as the ordered triple
\begin{equation}
T_{\beta}:=(\beta_c,\beta,1).
\label{eq:Tbeta}
\end{equation}
The reciprocal triangle has horizontal leg $\alpha$, vertical leg $\delta$,
and hypotenuse $\alpha\delta$:
\begin{equation}
T_{\rho}:=(\alpha,\delta,\alpha\delta).
\label{eq:Trho}
\end{equation}
Define
\begin{equation}
\kappa:=\alpha\delta
====================

\frac{1}{\beta\beta_c}.
\label{eq:kappa}
\end{equation}
Then
\begin{align}
\kappa T_{\beta}
&=
\left(
\frac{\beta_c}{\beta\beta_c},
\frac{\beta}{\beta\beta_c},
\frac{1}{\beta\beta_c}
\right)\
&=
\left(
\frac{1}{\beta},
\frac{1}{\beta_c},
\frac{1}{\beta\beta_c}
\right)\
&=
(\alpha,\delta,\alpha\delta)\
&=T_{\rho}.
\label{eq:trianglesimilarity}
\end{align}
Therefore
\begin{equation}
\boxed{T_{\rho}=\kappa T_{\beta}.}
\label{eq:boxsimilarity}
\end{equation}
The two triangles are similar by a uniform positive scale factor. Their
orientation is identical because
\begin{equation}
\frac{\delta}{\alpha}
=====================

# \frac{\beta}{\beta_c}

\tan\theta,
\label{eq:sameslope}
\end{equation}
and correspondingly
\begin{equation}
\frac{\alpha}{\delta}
=====================

# \frac{\beta_c}{\beta}

\cot\theta.
\label{eq:samecot}
\end{equation}

This is stronger than the observation that reciprocal identities happen to
close algebraically. The reciprocal construction preserves the angle of the
original triangle while changing its scale. In the language used later in
this paper: it is the same relation in a different body.

\section{Why the old $\sin\alpha=1/\alpha$ line is wrong, and why it refuses to die}

The expression
\begin{equation}
\sin\alpha=\frac{1}{\alpha}
\label{eq:wrong}
\end{equation}
is not an identity of ordinary trigonometry when $\alpha$ is a real number
interpreted as a circular angle in radians. With the definition
$\alpha=1/\beta$, it is generically false.

The correct statement is instead
\begin{equation}
\alpha=\csc\theta
\end{equation}
and therefore
\begin{equation}
\sin\theta=\frac{1}{\alpha}.
\label{eq:correctsine}
\end{equation}
The persistent attraction of \eqref{eq:wrong} comes from suppressing the
distinction between the underlying state and the coordinate used to name it.

To make that distinction explicit, let $\Aspace$ denote an abstract
one-dimensional state manifold for the positive one-axis relativistic
relation. A point
\begin{equation}
a\in\Aspace
\end{equation}
represents the physical state without yet choosing a coordinate. Introduce
the following coordinate maps:
\begin{align}
q_{\beta}(a)&=\beta\in(0,1),\
q_{\theta}(a)&=\theta\in(0,\pi/2),\
q_{\chi}(a)&=\chi\in(0,\infty),\
q_{\alpha}(a)&=\alpha\in(1,\infty),\
q_{\delta}(a)&=\delta\in(1,\infty).
\label{eq:charts}
\end{align}
They are related by
\begin{align}
q_{\beta}(a)
&=
\sin(q_{\theta}(a))
===================

# \tanh(q_{\chi}(a))

\frac{1}{q_{\alpha}(a)},
\label{eq:chartbeta}\
\sqrt{1-q_{\beta}(a)^2}
&=
\cos(q_{\theta}(a))
===================

# \sech(q_{\chi}(a))

\frac{1}{q_{\delta}(a)}.
\label{eq:chartbetac}
\end{align}

Now define the scalar function
\begin{equation}
S:\Aspace\rightarrow(0,1)
\end{equation}
by
\begin{equation}
S(a):=\sin(q_{\theta}(a)).
\label{eq:intrinsicsine}
\end{equation}
Then
\begin{equation}
S(a)
====

# q_{\beta}(a)

# \tanh(q_{\chi}(a))

\frac{1}{q_{\alpha}(a)}.
\label{eq:sameS}
\end{equation}
Nothing has happened to the ordinary sine function. It still receives an
ordinary circular angle $q_{\theta}(a)$ as its argument. What has changed is
that the physical object $a$ is no longer identified with any one of its
scalar coordinates.

This gives the title its precise meaning. The slogan
\begin{equation}
\texttt{sin()}\text{ stays }\texttt{sin()}
\end{equation}
means that the circular sine function is not deformed. The phrase
\begin{equation}
\text{angles become plural}
\end{equation}
means that the object whose orientation is being represented is permitted
more than one angular or angle-like coordinate body.

If one writes $\sin\alpha$ as shorthand for ``evaluate the circular sine
after converting the reciprocal coordinate $\alpha$ to the circular chart,''
then one is introducing an implicit coercion map
\begin{equation}
C_{\alpha\rightarrow\theta}(\alpha)
:=
\arcsin\left(\frac{1}{\alpha}\right).
\label{eq:coercion}
\end{equation}
The type-correct expression is
\begin{equation}
\sin\left(C_{\alpha\rightarrow\theta}(\alpha)\right)
====================================================

\frac{1}{\alpha}.
\label{eq:typecorrect}
\end{equation}
The provocative expression \eqref{eq:wrong} is obtained only by suppressing
$C_{\alpha\rightarrow\theta}$.

The missing map is the hole in the notation.

\section{Is $(\alpha,\delta)$ actually a chart?}

A point requiring more care than the terse version allowed is that
$\Aspace$ is one-dimensional. The pair $(\alpha,\delta)$ therefore cannot be
two independent coordinates on the one-axis state. Equation
\eqref{eq:reciprocalconstraint} constrains the pair to the curve
\begin{equation}
\mathcal{C}*{\rho}
:=
\left{
(\alpha,\delta)\in(1,\infty)^2:
\frac{1}{\alpha^2}+\frac{1}{\delta^2}=1
\right}.
\label{eq:reciprocalcurve}
\end{equation}
It is cleaner to regard either $\alpha$ or $\delta$ as a scalar coordinate
and the ordered pair as a redundant embedding
\begin{equation}
R:\Aspace\rightarrow\mathcal{C}*{\rho}\subset\mathbb{R}^2,
\qquad
R(a)=\bigl(q_{\alpha}(a),q_{\delta}(a)\bigr).
\label{eq:reciprocalembedding}
\end{equation}
The two entries make the complementary structure visible, but they do not add
a degree of freedom.

This distinction matters for any later R7 extension. If $\alpha$ and
$\delta$ acquire independent dynamical content in a higher-dimensional
model, that independence must arise from additional degrees of freedom in
the lifted geometry. It cannot be inferred from the one-axis reciprocal
triangle alone.

\section{Temporal and spatial vantage: what the one-axis mathematics supports}

Earlier GPD language associated $\alpha$ and $\delta$ with complementary
temporal and spatial aspects of a vantage tilt. The one-axis mathematics
supports a precise but limited version of that intuition.

Define
\begin{equation}
\beta_v:=\frac{v}{c},
\qquad
\beta_c:=\sqrt{1-\beta_v^2}.
\label{eq:betavbetac}
\end{equation}
Then
\begin{equation}
\alpha:=\frac{1}{\beta_v},
\qquad
\delta:=\frac{1}{\beta_c}.
\label{eq:alphadeltavc}
\end{equation}
If one defines
\begin{equation}
\gamma_v:=\frac{1}{\sqrt{1-\beta_v^2}},
\qquad
\gamma_c:=\frac{1}{\sqrt{1-\beta_c^2}},
\label{eq:twogammas}
\end{equation}
then, on the positive branch,
\begin{equation}
\boxed{
\delta=\gamma_v,
\qquad
\alpha=\gamma_c.
}
\label{eq:alphadeltagamma}
\end{equation}
The subscripts are easy to misread. The reciprocal of the velocity fraction
is the Lorentz-like factor built from the complementary fraction, while the
reciprocal of the complementary fraction is the ordinary Lorentz factor built
from the velocity fraction.

Explicitly,
\begin{align}
\gamma_c
&=
\frac{1}{\sqrt{1-\beta_c^2}}
============================

# \frac{1}{\sqrt{\beta_v^2}}

\frac{1}{\beta_v}
=\alpha,
\label{eq:gammacalpha}\
\gamma_v
&=
\frac{1}{\sqrt{1-\beta_v^2}}
============================

\frac{1}{\beta_c}
=\delta.
\label{eq:gammavdelta}
\end{align}
This crossed pairing is not an additional physical law. It is an algebraic
consequence of complementarity. It is nevertheless exactly the feature that
makes the reciprocal notation interesting: each reciprocal aperture is
naturally expressed as the gamma factor of the complementary coordinate.

In the one-axis case, therefore,
\begin{equation}
(\alpha,\delta)
===============

# (\gamma_c,\gamma_v)

# (\csc\theta,\sec\theta)

(\coth\chi,\cosh\chi).
\label{eq:masterpair}
\end{equation}

\section{The componentwise $c_i$ question}

The notation becomes substantially more delicate when one tries to promote
the one-axis complement to several spatial axes.

For a single selected axis $i$, one may define an algebraic complementary
quantity $c_i$ by
\begin{equation}
c^2=v_i^2+c_i^2,
\label{eq:componentcdef}
\end{equation}
so that
\begin{equation}
\beta_i:=\frac{v_i}{c},
\qquad
\beta_{c_i}:=\frac{c_i}{c}
\label{eq:componentbetas}
\end{equation}
and hence
\begin{equation}
\beta_i^2+\beta_{c_i}^2=1.
\label{eq:componentcircle}
\end{equation}
Axis by axis, the entire one-dimensional construction can then be repeated:
\begin{align}
\beta_i&=\sin\theta_i=\tanh\chi_i,
\label{eq:componentbeta}\
\beta_{c_i}&=\cos\theta_i=\sech\chi_i,
\label{eq:componentbetac}\
\alpha_i&:=\frac{1}{\beta_i}
=\csc\theta_i
=\coth\chi_i,
\label{eq:componentalpha}\
\delta_i&:=\frac{1}{\beta_{c_i}}
=\sec\theta_i
=\cosh\chi_i.
\label{eq:componentdelta}
\end{align}
One may likewise define
\begin{equation}
\gamma_i:=\frac{1}{\sqrt{1-\beta_i^2}},
\qquad
\gamma_{c_i}:=\frac{1}{\sqrt{1-\beta_{c_i}^2}},
\label{eq:componentgammas}
\end{equation}
which gives
\begin{equation}
\delta_i=\gamma_i,
\qquad
\alpha_i=\gamma_{c_i}
\label{eq:componentcrosspair}
\end{equation}
on the positive branch.

This is mathematically consistent axis by axis. It is not yet a consistent
three-dimensional relativistic interpretation.

For an ordinary three-velocity
\begin{equation}
\boldsymbol{\beta}
==================

(\beta_x,\beta_y,\beta_z),
\qquad
\beta^2
=======

\beta_x^2+\beta_y^2+\beta_z^2,
\label{eq:globalbeta}
\end{equation}
the physical Lorentz factor is
\begin{equation}
\gamma
======

# \frac{1}{\sqrt{1-\beta^2}}

\frac{1}{\sqrt{1-\beta_x^2-\beta_y^2-\beta_z^2}}.
\label{eq:globalgamma}
\end{equation}
By contrast, the axiswise quantity
\begin{equation}
\delta_i
========

\frac{1}{\sqrt{1-\beta_i^2}}
\label{eq:axisdelta}
\end{equation}
is generally not equal to $\gamma$ when more than one component of velocity
is nonzero.

There are therefore two distinct constructions that must not be conflated.

First, a \emph{global complement} may be defined by
\begin{equation}
\beta_c
:=
\sqrt{1-\lVert\boldsymbol{\beta}\rVert^2},
\label{eq:globalcomplement}
\end{equation}
so that
\begin{equation}
\beta^2+\beta_c^2=1,
\qquad
\gamma=\beta_c^{-1}.
\label{eq:globalcomplementrelation}
\end{equation}
This is directly tied to ordinary special relativity.

Second, one may define a family of \emph{axiswise complements}
\begin{equation}
\beta_{c_i}
:=
\sqrt{1-\beta_i^2},
\qquad
c_i:=c,\beta_{c_i}.
\label{eq:axiscomplements}
\end{equation}
Each pair $(\beta_i,\beta_{c_i})$ lies on its own unit circle, but the $c_i$
are not thereby established as Cartesian components of one conventional
velocity vector. They are three separate complements, each defined relative
to a selected axis.

This is precisely where the R7 proposal becomes relevant rather than
decorative. If the lifted geometry supplies a distinct conjugate or time-like
partner for each spatial axis, then the three quantities $c_i$ may acquire a
natural home as components in that larger structure. Without such a
construction, calling them ``components of $c$'' would overstate what the
one-axis algebra proves.

Accordingly, the componentwise hypothesis
\begin{equation}
\boxed{
\beta_{c_i}=\frac{c_i}{c}
}
\label{eq:componenthypothesis}
\end{equation}
should be retained as an explicit extension problem, not silently imported
into the one-axis result.

\section{The status of $\varphi$}

The terse draft moved too quickly over the third symbol $\varphi$ that
appeared in the earlier GPD triangle construction.

If $\alpha$ and $\delta$ are ordinary Euclidean angles of one triangle, then
expressions involving
\begin{equation}
\pi-(\alpha+\delta)
\label{eq:ordinaryphi}
\end{equation}
are meaningful in the usual way. But the present reinterpretation assigns
\begin{equation}
\alpha=\csc\theta,
\qquad
\delta=\sec\theta,
\label{eq:alphadeltaagain}
\end{equation}
so $\alpha$ and $\delta$ are reciprocal coordinates, not circular angles
measured in radians. Under that interpretation, the sum $\alpha+\delta$
cannot simply be subtracted from $\pi$ and called an angular remainder.
Doing so mixes coordinate types.

At the same time, the reciprocal pair contains a natural mixed ratio:
\begin{equation}
\frac{\delta}{\alpha}
=====================

# \tan\theta

\sinh\chi.
\label{eq:mixedratio}
\end{equation}
This quantity relates the two reciprocal apertures and reconstructs either
the circular or hyperbolic parameter:
\begin{equation}
\theta
======

\arctan\left(\frac{\delta}{\alpha}\right),
\label{eq:thetarecover}
\end{equation}
\begin{equation}
\chi
====

\operatorname{arsinh}\left(\frac{\delta}{\alpha}\right).
\label{eq:chirecover}
\end{equation}

It is tempting to identify the old $\varphi$ with this transition relation.
That identification is not yet justified. The safer conclusion is narrower:
once $\alpha$ and $\delta$ are retyped as reciprocal coordinates, the old
Euclidean interpretation of $\varphi$ must be re-derived. In a genuinely
higher-dimensional vantage geometry, $\varphi$ may encode an independent
mixed spatial--temporal orientation; in the one-axis construction, no
independent third degree of freedom has yet appeared.

Thus $\varphi$ remains an open variable, not a casualty and not a solved
quantity.

\section{Why this helps motivate lift--rotate--squeeze}

A higher-dimensional Euclidean reconstruction of special relativity faces a
severe methodological objection. If one begins with the Lorentz
transformation, invents a lift, chooses a rotation, and then chooses a squeeze
specifically so that the final expression reproduces the Lorentz
transformation, the construction may be algebraically correct while
explaining nothing. It is a factorization, not a derivation.

The present observation does not remove that burden. It does, however,
supply an independent reason to expect both a circular rotation and a
nontrivial rescaling to appear.

The circular coordinate $\theta$ and rapidity $\chi$ describe the same
scalar speed parameter through
\begin{equation}
\beta=\sin\theta=\tanh\chi.
\label{eq:twoparameterizations}
\end{equation}
The Gudermannian provides the exact real transition map
\begin{equation}
\theta=\gd\chi.
\label{eq:gdtransition}
\end{equation}
The reciprocal representation then produces a second triangle related to the
normalized circular triangle by the scale factor
\begin{equation}
\kappa=\alpha\delta.
\label{eq:kappaagain}
\end{equation}
Thus the elementary scalar geometry itself contains both
\begin{equation}
\text{an orientation variable }\theta
\end{equation}
and
\begin{equation}
\text{a reciprocal scale }\kappa.
\end{equation}
A lift--rotate--squeeze construction can therefore be motivated schematically
as an attempt to realize these two features geometrically:
\begin{equation}
\text{state}
\xrightarrow{\text{lift}}
\text{vantage geometry}
\xrightarrow{\text{rotate by }\theta}
\text{reoriented state}
\xrightarrow{\text{squeeze}}
\text{reported relativistic coordinates}.
\label{eq:lrs}
\end{equation}

The crucial word is \emph{attempt}. The scalar identities do not prove that
a Euclidean rotation in R7 is dynamically or geometrically equivalent to a
Lorentz boost. In particular, a successful construction must still answer at
least the following questions.

\begin{enumerate}[label=(\arabic*)]
\item What is the precise lifted manifold or vector space?
\item What metric, bilinear form, or other invariant structure does it carry?
\item Which subspace is rotated, and why is $\theta$ the rotation parameter
there?
\item What is the squeeze map, and is it uniquely determined by the geometry
rather than fitted to the Lorentz transformation?
\item How does the construction compose under successive boosts?
\item How are non-collinear boosts, Thomas--Wigner rotation, and the global
velocity norm recovered?
\item What role, if any, do the axiswise complements $c_i$ play in the lifted
space?
\end{enumerate}

The conceptual gain is therefore modest but real: circular rotation is no
longer introduced merely because Euclidean rotations are aesthetically
attractive. A standard real map already connects the circular angle to
rapidity, and the reciprocal construction already supplies a natural scaling
structure. The higher-dimensional theory must explain why those scalar facts
lift to the proposed geometry.

\section{What ``angles become plural'' should mean}

The phrase can be stated weakly or strongly.

The weak statement is mathematically uncontroversial. A single
one-parameter relativistic state admits several smooth coordinates related by
invertible maps on the chosen branch:
\begin{equation}
\beta
\longleftrightarrow
\theta
\longleftrightarrow
\chi
\longleftrightarrow
\alpha
\longleftrightarrow
\delta.
\label{eq:coordinatechain}
\end{equation}
Calling all of them ``angles'' would be nonstandard terminology, but the
coordinate equivalence itself is elementary.

The stronger proposal is geometric. It asserts that the underlying object
should be modeled so that these coordinates are not merely substitutions
made after the fact, but natural projections or charts selected by different
vantages. In such a theory, ``angle'' would mean something closer to an
orientation relation or group-theoretic state than to a single real number
measured as Euclidean arc length divided by radius.

This stronger statement is compatible with familiar mathematical practice
in spirit: an angle need not be identified ontologically with one chosen
coordinate. Circular rotations are elements of $SO(2)$, while collinear
Lorentz boosts form a one-parameter subgroup whose additive parameter is
rapidity. Coordinates label those transformations; they are not the
transformations themselves.

The present proposal asks whether GPD's vantage object can be defined at that
level, with
\begin{equation}
\theta,
\qquad
\chi,
\qquad
\alpha,
\qquad
\delta
\end{equation}
serving as different readable bodies of one relation.

That is the sense in which angles become plural.

\section{Consolidated identity table}

For $0<\beta<1$, define
\begin{equation}
\beta_c=\sqrt{1-\beta^2},
\qquad
\theta=\arcsin\beta,
\qquad
\chi=\artanh\beta,
\qquad
\alpha=\frac{1}{\beta},
\qquad
\delta=\frac{1}{\beta_c}.
\label{eq:definitionsummary}
\end{equation}
Then
\begin{align}
\beta
&=\sin\theta
=\tanh\chi
=\frac{1}{\alpha},
\label{eq:summary1}\
\beta_c
&=\cos\theta
=\sech\chi
=\frac{1}{\delta},
\label{eq:summary2}\
\alpha
&=\csc\theta
=\coth\chi,
\label{eq:summary3}\
\delta
&=\sec\theta
=\cosh\chi,
\label{eq:summary4}\
\frac{\delta}{\alpha}
&=\tan\theta
=\sinh\chi,
\label{eq:summary5}\
\frac{\alpha}{\delta}
&=\cot\theta
=\csch\chi,
\label{eq:summary6}\
\frac{1}{\alpha^2}+\frac{1}{\delta^2}
&=1,
\label{eq:summary7}\
\alpha^2+\delta^2
&=\alpha^2\delta^2,
\label{eq:summary8}\
(\alpha,\delta,\alpha\delta)
&=
\alpha\delta,(\beta_c,\beta,1).
\label{eq:summary9}
\end{align}
If
\begin{equation}
\gamma_v:=\frac{1}{\sqrt{1-\beta^2}},
\qquad
\gamma_c:=\frac{1}{\sqrt{1-\beta_c^2}},
\end{equation}
then
\begin{equation}
\boxed{
\delta=\gamma_v,
\qquad
\alpha=\gamma_c.
}
\label{eq:summarygamma}
\end{equation}
Every equality in this section follows from the definitions and standard
identities. None requires GPD.

\section{The proposed promotion}

The speculative content can now be stated without borrowing any of the
algebra's authority.

Let $a\in\Aspace$ denote a vantage relation. The proposal is to investigate a
geometry in which
\begin{equation}
q_{\theta}(a)=\theta,
\qquad
q_{\chi}(a)=\chi,
\qquad
q_{\alpha}(a)=\alpha,
\qquad
q_{\delta}(a)=\delta
\end{equation}
are natural representations of that same relation, with transition maps
fixed by the identities above. Circular angle, hyperbolic rapidity, and
reciprocal aperture would then be interpreted not as rival descriptions but
as different coordinate bodies exposed by different projections of the same
object.

Under that promotion, the old notational crime
\begin{equation}
\sin\alpha=\frac{1}{\alpha}
\end{equation}
can be read as evidence of a suppressed chart transition rather than as a new
trigonometric identity. The type-correct statement remains
\begin{equation}
\sin\left(\arcsin\frac{1}{\alpha}\right)
========================================

\frac{1}{\alpha}.
\end{equation}
The function is untouched. The argument has changed bodies.

The research question is whether a lifted physical geometry can make that
change of body structural rather than notational.

\section{Conclusion}

The one-axis mathematics is small.

A relativistic speed parameter admits a circular representation and a
hyperbolic representation connected by the Gudermannian. The reciprocal
quantities that appeared in earlier GPD sketches are exactly cosecant/secant
in the circular representation and coth/cosh in the hyperbolic
representation. Their constrained pair forms a reciprocal triangle that is
similar to the original normalized triangle by the scale factor
$\alpha\delta$. The crossed identities $\alpha=\gamma_c$ and
$\delta=\gamma_v$ follow directly from the complementary definitions.

The interpretation is larger.

Instead of modifying trigonometric functions to accommodate the reciprocal
notation, one may keep the functions ordinary and promote the represented
relation. An angle, in this sense, is not one privileged scalar. It is an
orientation-like object admitting multiple coordinate bodies. The circular
angle $\theta$, rapidity $\chi$, and reciprocal apertures $(\alpha,\delta)$
become different ways of reading the same one-parameter relation.

This does not establish the GPD R7 construction. It does sharpen what that
construction would have to accomplish. In particular, the componentwise
proposal $\beta_{c_i}=c_i/c$ must be distinguished from the global
relativistic complement, and any lift--rotate--squeeze derivation must explain
rather than fit its lift, rotation, and squeeze.

The seed is therefore deliberately narrower than the surrounding theory:
\begin{equation}
\boxed{
\texttt{sin()}\text{ stays }\texttt{sin()};\quad
\text{angles become plural.}
}
\end{equation}

\begin{center}
\emph{Same relation. Different body.}
\end{center}

\begin{thebibliography}{9}

\bibitem{MunroGPD}
W. B. Munro,
\emph{General Peace Dynamics},
working manuscript and special-relativity sketch notes, 2025.

\bibitem{DLMF}
NIST Digital Library of Mathematical Functions,
\S 4.23(viii), ``Gudermannian Function,''
Eqs. 4.23.39--4.23.42.

\end{thebibliography}

\end{document}
